% Generated draft: requires mathematical and source review.
\documentclass[11pt]{article}
\usepackage[margin=25mm]{geometry}
\usepackage{fontspec}
\setmainfont{Latin Modern Roman}
\usepackage{amsmath,amssymb,mathtools,unicode-math}
\setmathfont{Latin Modern Math}
\usepackage{enumitem,needspace,xurl,hyperref}
\setlength{\parindent}{0pt}
\setlength{\parskip}{5pt}
\setlength{\emergencystretch}{3em}
\newcommand{\N}{\mathbb{N}}
\newcommand{\Z}{\mathbb{Z}}
\newcommand{\Q}{\mathbb{Q}}
\newcommand{\R}{\mathbb{R}}
\newcommand{\C}{\mathbb{C}}
\newcommand{\F}{\mathbb{F}}
\newcommand{\A}{\mathbb{A}}
\newcommand{\PP}{\mathbb{P}}
\newcommand{\E}{\mathbb{E}}
\newcommand{\eps}{\varepsilon}
\newcommand{\dd}{\,\mathrm{d}}
\begin{document}
\section*{UnsolvedMath ID: 6700027}
\textbf{Source problem code:} catalogue entry
\subsection{Definitions and mathematical statement}
Let $(X,g)$ be a Riemannian manifold and let $\delta>0$. For a surface $S$ in $X$, its metric $\delta$-neighborhood is
\[
 U_\delta(S)=\{x\in X:\operatorname{dist}_g(x,S)<\delta\}.
\]
The source singles out ``$T(X)$-non-spin surfaces.'' This refers to the spin behavior of the ambient tangent bundle along a surface, not automatically to the surface's own tangent bundle. In an oriented pullback-bundle interpretation, one would test the second Stiefel--Whitney class of the restricted or pulled-back $TX$, but the source gives no exact surface or bundle convention.

The proposed qualitative implication is that if the neighborhoods $U_\delta(S)$ of every such surface are ``large,'' then the spin area of $X$ must also be large. The supplied record defines neither spin area nor largeness. In particular it does not specify whether neighborhood size means volume, diameter, width, a surface-filling quantity, or another invariant; nor does it prescribe any lower bound or dependence on $\delta$.

A complete quantitative statement would require an admissible surface class, the neighborhood-size functional, a precise spin-area invariant, and the inequality linking their values. It must also address the case of no qualifying non-spin surfaces and specify the dimension, compactness, and orientation assumptions on $X$. These choices cannot be recovered from the qualitative sentence and are not silently supplied.
\subsection{Short English statement}
Does uniformly large neighborhood geometry around all ambient-tangent-non-spin surfaces force large spin area of the manifold? The relevant invariants and thresholds are missing, so the source gives a speculative program rather than a quantified inequality.

The supplied summary connects the idea with spin area, $K$-area, slice-area, and two-dimensional scalar-curvature obstructions, but says no direct theorem verifying the stated implication was found. It gives no definition equating those notions. The OPEN-TRIAGE label and 2017 passage are retained with current status needing review.
\subsection{Sources}
\begin{itemize}
\item Local UnsolvedMath registry, permanent ID 6700027, source code catalogue entry, scalar-curvature Question [?27], proposed by Misha Gromov, 2017. Dataset curation and metadata: Ulam AI and UnsolvedMath Contributors, CC BY 4.0.
\item Misha Gromov, \emph{101 Questions, Problems and Conjectures around Scalar Curvature} (2017), printed tag [?27], occurrence 28, PDF page 38, \url{https://www.ihes.fr/~gromov/wp-content/uploads/2018/08/101-problemsOct1-2017.pdf}; normalized pdf-text extraction, accessed 2026-07-29.
\item No exact later theorem or independent definition of spin area is supplied. Source-specific provenance: the author-hosted PDF is incomplete and publicly available; current status and statement-reuse terms NEEDS\_REVIEW, separately from the dataset licence.
\end{itemize}
\end{document}
